\documentclass{article}

\usepackage{amsfonts}
\usepackage{amsthm}
\usepackage{mathrsfs}
\usepackage{amsmath}
\usepackage{graphicx}
\usepackage{hyperref}
\usepackage{stackengine}
\newtheorem{theorem}{Theorem}
\theoremstyle{definition}\newtheorem{definition}{DEF}
\theoremstyle{lemma}\newtheorem{lemma}{Lemma}
\theoremstyle{corrolary}\newtheorem{corrolary}{Cor}
\setlength{\parindent}{0pt}

\begin{document}

\begin{definition}
A parametric distribution is a class of probability distributions indexed by finitely many variables.
\end{definition}

\begin{definition}
Let $a,b\in \mathbb{R}$.
Say $X$ has a uniform distribution, $X \sim Unif(a,b)$ when the PDF is $f(x)=\frac{1}{b-a}$ for $x\in (a,b)$.
\end{definition}

\begin{theorem}
THe CDF is $F(x)=\frac{x-a}{b-a}$ and MGF $M(t)=\frac{e^{tb}-e^{ta}}{t(b-a)}$.
\end{theorem}
\begin{proof}
Assume $a<x<b$.
For the CDF,
\begin{align*}
F(x)&=\int_a^x 1dF\\
&=\int_a^x \frac{1}{b-a}dy\\
&=\frac{y}{b-a}|_a^x\\
&=\frac{x-a}{b-a}
\end{align*}
For the MGF,
\begin{align*}
M(t)&=\mathbb{E}(e^{tX})\\
&= \int_a^b e^{tx}dF\\
&=\int_a^b \frac{e^{tx}}{b-a}dx\\
&=\frac{e^{tx}}{t(b-a)}|_a^b\\
&=\frac{e^{tb}-e^{ta}}{t(b-a)}
\end{align*}
\end{proof}
\begin{corrolary}
$\mathbb{E}(X)=\frac{a+b}{2}$ and $Var(X)=\frac{(b-a)^2}{12}$.
\end{corrolary}

\begin{definition}
A Bernouli trial is a random experiment with two outcomes: success, which happens with probability $p$, and failure, which happens with probability $1-p$.
\end{definition}

\begin{definition}
Let $p\in (0,1)$.
Say a random variable has a Bernoouli distribution whenever it has pdf $f(k)=p^k(1-p)^{1-k}$ for $k\in \{0,1\}$.
\end{definition}

\begin{theorem}
A Bernouli distribution has CDF $F(k)=1-p+kp$ and MGF $M(t)=1-p+pe^t$.
\end{theorem}
\begin{proof}
For the CDF,
\begin{align*}
F9x)&=\int_0^k1dF\\
&=\sum_{i=0}^k p^i(1-p)^{1-i}\\
&=1-p+kp
\end{align*}
Note that the easiest way to do this summation is to simply plug in both $k=0$ and $k=1$ and find the values.
For the MGF, we perform a similar process
\begin{align*}
M(t)&=\mathbb{E}(e^{tX})\\
&=\int_0^ke^{ti}dF\\
&=\sum_{i=0}^k e^{ti}p^i(1-p)^{1-i}\\
&=1-p+pe^t
\end{align*}
\end{proof}
\begin{corrolary}
$\mathbb{E}(X)=p$ and $Var(X)=p(1-p)$.
\end{corrolary}

\textbf{Ex-}
A coin has a $50\%$ chance of being hads and will either be hads or tails.
Thus, flipping a coin is a Bernouli trial, and whether it comes up heads follows a Bernouli distribution.

\begin{definition}
Let $p\in(0,1)$ and $n\in\mathbb{Z}^+$.
Say a random vairable has a binomial distribution, $X\sim Binom(p,n)$, whenever i has PDF $f(k)=\binom{n}{k} p^k (1-p)^{n-k}$ for $k\in \{0, 1, 2, ..., n\}$.
\end{definition}

\begin{theorem}
A binomial distribution has MGF $M(t)=\left( 1-p+pe^t\right)^n$.
\end{theorem}
\begin{proof}
Notice
\begin{align*}
M(t)&=\mathbb{E}(e^{tX})\\
&=\int_0^k e^{ti}dF\\
&=\sum_{i=0}^k e^{ti}\binom{n}{i}p^i(1-p)^{n-i}\\
&=\sum_{i=0}^k \binom{n}{i}\left( pe^t\right)^i(1-p)^{n-i}\\
&=\left(1-p+pe^t\right)^n
\end{align*}
For the final step, we have used the binomial theorem.
\end{proof}
\begin{corrolary}
$\mathbb{E}(X)=np$ and $Var(X)=np(1-p)$.
\end{corrolary}

\textbf{Note-} The binomail distribution does not have a closed-form representation of the CDF, so we have not definte one here.
The CDF can be defined as a simple summation, but that sum does not have a simple evaluation in general.\\

\textbf{Ex-}
If a fair coin is flipped 10 times, the results are independent.
Thus, the number of heads resulting from 10 cooin flips follows a binomial distribution.

\begin{theorem}
Let $p\in (0,1)$, $n\in \mathbb{Z}^+$, and $X_k  { {iid} \atop \sim} Bernouli(p)$ for $1\leq k\leq n$.
Then
\begin{equation}
\sum_{k=1}^n X_k \sim Binom(p,n)
\end{equation}
\end{theorem}
\begin{proof}
Let $M_X$ be the MGF of $X_k$.
Notice the MGF is the same for all $k$ because the $X_k$ are defined to be iid.
Thus, the MGF of the sum is
\begin{align*}
M(t)&=\left(M(t)\right)^n\\
&=\left( 1-p+pe^t\right)^n
\end{align*}
This is the MGF of the binomial distribution.
Thus, the sum follows the binomial distribution.
\end{proof}

\textbf{Note-}
The binomial distribution can be described as the probability that $n$ independent Bernouli trials result in $k$ successes.
This situation often arises in applications of the probability distribution.

\begin{definition}
Let $p\in (0,1)$.
Say a random variable has a geometric distribution, $X\sim Geo(p)$, whenever it has PDF $f(k)=p(1-p)^k$ for $k\in \mathbb{Z}^+$.
\end{definition}

\begin{theorem}
The geometric distribution has CDF $F(k)=1-(1-p)^k$ and MGF $M(t)=\frac{p}{1-(1-p)e^t}$.
\end{theorem}
\begin{proof}
In order to compute the CDF, we consider
\begin{align*}
Pr(X>k)&=\int_k^\infty 1dF\\
&=\sum_k^\infty p(1-p)^i\\
&=(1-p)^k
\end{align*}
This summation is an example of a geometric series, which is where the distribution gets its name.
To compute the CDF, we simply take the compliment of htis event, so $F(k)=1-p(1-p)^k$.\\
To compute the MGF, we can compute a similar geometric series
\begin{align*}
M(t)&=\mathbb{E}(e^{tX})\\
&=\int_0^\infty 1dF\\
&=\sum_{k=0}^\infty e^{tk}p(1-p)^k\\
&=\sum_{k=0}^\infty p(e^t(1-p))^k\\
&=\frac{p}{1-(1-p)e^t}
\end{align*}
\end{proof}
\begin{corrolary}
$\mathbb{E}(X)=\frac{1}{p}$ and $Var(X)=\frac{1-p}{p^2}$
\end{corrolary}

\textbf{Note-}
The geometric distribution can be interpreted as the probability that the first success in a series of independent Bernouli trials occurs on trial $k$.

\end{document}
